\noindent\textbf{Operaciones y resultados intermedios.}

\begin{enumerate}

\item Área algebraica de la sección, descontando los huecos.

\[\fitmath{y_{1}=\left(80-0\right)\cdot \left(160-0\right)\approx 12800\,\mathrm{mm^{2}}}\]

\item Centroide en el eje vertical mediante primeros momentos.

\[\fitmath{y_{2}=\frac{\frac{\left(80-0\right)\cdot \left(160-0\right)\cdot \left(0+160\right)}{2}}{\left(80-0\right)\cdot \left(160-0\right)}\approx 80\,\mathrm{mm}}\]

\item Momento de inercia I\_z mediante el teorema de Steiner.

\[\fitmath{y_{3}=\frac{\left(80-0\right)\cdot {\left(160-0\right)}^{3}}{12}+\left(80-0\right)\cdot \left(160-0\right)\cdot {\left(\frac{0+160}{2}-80\right)}^{2}\approx 2{,}730667\times10^{7}\,\mathrm{mm^{4}}}\]

\item Tensión extrema: axil y flexión con el convenio declarado.

\[\fitmath{y_{4}=\frac{-525}{32}\approx -16{,}40625\,\mathrm{MPa}}\]

\end{enumerate}
